\begin{PSbox}[unbreakable]{obj}{Learning Objectives}

After completing this module, students will be able to:

\begin{enumerate}
\def\labelenumi{\arabic{enumi}.}
\tightlist
\item
  Distinguish between population and sample
\item
  Distinguish between parameter and statistic
\item
  Define and compute sample mean and sample variance
\item
  Explain sampling distributions and why they matter
\item
  Understand the transition from probability $(\text{model}\ \to \text{predictions})$ to statistics $(\text{data}\ \to \text{model})$
\end{enumerate}

\end{PSbox}

\PSsection{12.1}{Introduction: From Probability to Statistics}

\PSsubsection{The Two Directions}

{\def\LTcaptype{none} % do not increment counter
\begin{longtable}[c]{@{}cc@{}}
\toprule\noalign{}
\rowcolor{PSnavy}\PSth{Probability} & \PSth{Statistics} \\
\midrule\noalign{}
\endhead
\bottomrule\noalign{}
\endlastfoot
Model is KNOWN & Model is UNKNOWN \\
Predict outcomes & Infer model from data \\
Forward problem & Inverse problem \\
``Given distribution, what data?'' & ``Given data, what distribution?'' \\
\end{longtable}
}

\textbf{Probability} (Modules 1–11): Given a model \ensuremath{\to} predict observations.\PSbr{}\textbf{Statistics} (Modules 12–16): Given observations \ensuremath{\to} infer the model.

\PSsubsection{Engineering Context}

\begin{itemize}
\tightlist
\item
  \textbf{Probability:} ``If noise is $N(0,\, 0.01)$, what is $P(\text{error})$?''
\item
  \textbf{Statistics:} ``I measured 1000 noise samples. What is $\sigma^{2}$?''
\end{itemize}

\PSsection{12.2}{Population vs. Sample}

\PSsubsection{Population}

The \textbf{complete} set of ALL items/measurements of interest.

{\def\LTcaptype{none} % do not increment counter
\begin{longtable}[c]{@{}cc@{}}
\toprule\noalign{}
\rowcolor{PSnavy}\PSth{Engineering Context} & \PSth{Population} \\
\midrule\noalign{}
\endhead
\bottomrule\noalign{}
\endlastfoot
All chips ever produced by a fab & ALL chip delays \\
The noise process forever & ALL noise voltage values \\
All network packets ever sent & ALL packet delays \\
\end{longtable}
}

Populations are usually \textbf{infinite or impractically large} — we cannot measure everything.

\PSsubsection{Sample}

A \textbf{subset} of the population that we actually observe/measure.

{\def\LTcaptype{none} % do not increment counter
\begin{longtable}[c]{@{}cc@{}}
\toprule\noalign{}
\rowcolor{PSnavy}\PSth{Engineering Context} & \PSth{Sample} \\
\midrule\noalign{}
\endhead
\bottomrule\noalign{}
\endlastfoot
50 chips tested from production line & Measured delays \\
10000 noise voltage samples recorded & Recorded values \\
500 packet delays measured & Observed delays \\
\end{longtable}
}

\PSsubsection{Why We Sample}

\begin{itemize}
\tightlist
\item
  Cannot test every chip (destructive testing)
\item
  Cannot record noise forever (finite memory)
\item
  Cannot measure every packet (too many)
\end{itemize}

The goal: \textbf{Infer properties of the population from the sample.}

\PSsection{12.3}{Parameter vs. Statistic}

\PSsubsection{Parameter}

A \textbf{fixed} (but unknown) number that describes the population:

\begin{itemize}
\tightlist
\item
  $\mu$ = true population mean
\item
  $\sigma^{2}$ = true population variance
\item
  $p$ = true defect rate
\item
  $\lambda$ = true failure rate
\end{itemize}

Parameters are \textbf{constants} — they have ONE true value (we just don’t know it).

\PSsubsection{Statistic}

A \textbf{function of sample data} used to estimate or test parameters:

\begin{itemize}
\tightlist
\item
  $\bar{X}$ = sample mean (estimates $\mu$)
\item
  $S^{2}$ = sample variance (estimates $\sigma^{2}$)
\item
  $\hat{p}$ = sample proportion (estimates $p$)
\end{itemize}

\begin{PSquote}

\PSwarn{} \textbf{Key insight:} Statistics are RANDOM VARIABLES! Different samples give different values.

\end{PSquote}

\begin{PSbox}[unbreakable]{ex}{Example}

True noise power: $\sigma^{2} = 0.04\,\mathrm{V}^{2}$ (parameter — fixed, unknown)

Experiment 1: Collect 100 samples \ensuremath{\to} compute $S^{2} = 0.038$ (one realization)\PSbr{}Experiment 2: Collect 100 samples \ensuremath{\to} compute $S^{2} = 0.042$ (different realization)\PSbr{}Experiment 3: Collect 100 samples \ensuremath{\to} compute $S^{2} = 0.041$ (another realization)

Each $S^{2}$ is different because of randomness in sampling.

\end{PSbox}

\PSsection{12.4}{Random Sample}

\begin{PSbox}[unbreakable]{def}{Definition}

A \textbf{random sample} $X_{1},\, X_{2},\, \ldots,\, X_{n}$ is a set of $n$ independent, identically distributed $(i.i.d.)$ random variables from the population distribution.

\end{PSbox}

\PSsubsection{Requirements}

\begin{enumerate}
\def\labelenumi{\arabic{enumi}.}
\tightlist
\item
  \textbf{Independent:} Each observation doesn’t influence others
\item
  \textbf{Identically distributed:} All come from the same population
\item
  \textbf{Random selection:} No systematic bias in choosing which items to measure
\end{enumerate}

\PSsubsection{When i.i.d. Applies}

\PSyes{} Random selection of components from a large batch\PSbr{}\PSyes{} Independent measurements of a stable process\PSbr{}\PSyes{} Noise samples separated by more than the correlation time

\PSsubsection{When i.i.d. Fails}

\PSno{} Consecutive samples of a correlated signal\PSbr{}\PSno{} Measurements during a drifting process\PSbr{}\PSno{} Non-random selection (e.g., only testing ``suspicious'' components)

\PSsection{12.5}{Sampling Distributions}

\begin{PSbox}[unbreakable]{def}{Definition}

The \textbf{sampling distribution} of a statistic is the probability distribution of that statistic across all possible samples of size $n$.

\end{PSbox}

\PSsubsection{Why It Matters}

A statistic (like $\bar{X}$) is a random variable. Its distribution tells us:

\begin{itemize}
\tightlist
\item
  How much will it vary from sample to sample?
\item
  How close is it likely to be to the true parameter?
\item
  How should we quantify our uncertainty?
\end{itemize}

This is the foundation for confidence intervals and hypothesis testing.

\PSsection{12.6}{Sample Mean $\bar{X}$}

\begin{PSbox}[unbreakable]{def}{Definition}

\PSdisplay{\bar{X} = \frac{1}{n}\sum_{i=1}^n X_i}

\end{PSbox}

\PSsubsection{Properties}

\textbf{Mean of $\bar{X}$:}\PSdisplay{E[\bar{X}] = \mu}

$\bar{X}$ is an \textbf{unbiased} estimator of $\mu$ (on average, it equals the truth).

\textbf{Variance of $\bar{X}$:}\PSdisplay{\text{Var}(\bar{X}) = \frac{\sigma^2}{n}}

More measurements \ensuremath{\to} less variability \ensuremath{\to} more precise estimate.

\textbf{Standard Error:}\PSdisplay{\text{SE}(\bar{X}) = \frac{\sigma}{\sqrt{n}}}

\PSsubsection{Sampling Distribution of $\bar{X}$}

\begin{itemize}
\tightlist
\item
  If population is Gaussian: $\bar{X} \sim N\left(\mu,\, \frac{\sigma^{2}}{n}\right)$ \textbf{exactly} for any $n$
\item
  If population is NOT Gaussian: $\bar{X} \approx N\left(\mu,\, \frac{\sigma^{2}}{n}\right)$ \textbf{approximately} for large $n$ (by CLT)
\end{itemize}

\PSsection{12.7}{Sample Variance $S^{2}$}

\begin{PSbox}[unbreakable]{def}{Definition}

\PSdisplay{S^2 = \frac{1}{n-1}\sum_{i=1}^n (X_i - \bar{X})^2}

\end{PSbox}

\PSsubsection{Why $n-1$? (Bessel’s Correction)}

Using $n$ in the denominator gives a \textbf{biased} estimator (systematically underestimates $\sigma^{2}$):

$\displaystyle E\left[\frac{1}{n} \cdot \sum (X_{i} - \bar{X})^{2}\right] = \frac{n-1}{n}\sigma^{2} \ne \sigma^{2}$

Using $n-1$ corrects this:\PSdisplay{E[S^2] = \sigma^2}

\textbf{Intuition:} $\bar{X}$ is computed from the same data, so deviations from $\bar{X}$ are smaller than deviations from $\mu$. We ``lose one degree of freedom'' by estimating the mean.

\PSsubsection{Degrees of Freedom}

With $n$ data points and $\bar{X}$ estimated:

\begin{itemize}
\tightlist
\item
  The $n$ deviations $(X_{i} - \bar{X})$ sum to zero: $\sum (X_{i} - \bar{X}) = 0$
\item
  Only $n-1$ are ``free'' — the last is determined
\item
  $df = n - 1$
\end{itemize}

\PSsubsection{Sampling Distribution (Normal Population)}

If population is $N(\mu,\, \sigma^{2})$:\PSdisplay{(n-1)S^2/\sigma^2 \sim \chi^2(n-1)}

\PSsection{12.8}{Sample Standard Deviation}

\PSdisplay{S = \sqrt{S^2} = \sqrt{\frac{1}{n-1}\sum_{i=1}^n (X_i - \bar{X})^2}}

Note: $E[S] \ne \sigma$ (S is slightly biased for $\sigma$, though $S^{2}$ is unbiased for $\sigma^{2}$).

\PSsection{12.9}{Connection to CLT}

From Module 11: $\bar{X} \to N\left(\mu,\, \frac{\sigma^{2}}{n}\right)$ as $n \to \infty$.

This means the sampling distribution of $\bar{X}$ is approximately Gaussian. This is what allows us to:

\begin{enumerate}
\def\labelenumi{\arabic{enumi}.}
\tightlist
\item
  Build confidence intervals (Module 14)
\item
  Perform hypothesis tests (Module 15)
\item
  Quantify uncertainty in estimates
\end{enumerate}

\PSsection{12.10}{MATLAB Examples}

\begin{PSbox}[unbreakable]{ex}{Example 1: Sampling Distribution of $\bar{X}$}

\tcbset{PScodemode/.style={unbreakable}}

\begin{Shaded}
\begin{Highlighting}[]
\CommentTok{\%\% Demonstrate: X̄ is a random variable with its own distribution}
\CommentTok{\% Population: Exponential(λ=1), μ=1, σ²=1}
\VariableTok{mu\_pop} \OperatorTok{=} \FloatTok{1}\OperatorTok{;} \VariableTok{sigma2\_pop} \OperatorTok{=} \FloatTok{1}\OperatorTok{;}
\VariableTok{n} \OperatorTok{=} \FloatTok{25}\OperatorTok{;}  \CommentTok{\% Sample size}
\VariableTok{N\_experiments} \OperatorTok{=} \FloatTok{10000}\OperatorTok{;}

\CommentTok{\% Repeat experiment N\_experiments times}
\VariableTok{X\_bar} \OperatorTok{=} \VariableTok{zeros}\NormalTok{(}\FloatTok{1}\OperatorTok{,} \VariableTok{N\_experiments}\NormalTok{)}\OperatorTok{;}
\KeywordTok{for} \VariableTok{i} \OperatorTok{=} \FloatTok{1}\OperatorTok{:}\VariableTok{N\_experiments}
    \VariableTok{sample} \OperatorTok{=} \VariableTok{exprnd}\NormalTok{(}\FloatTok{1}\OperatorTok{,} \FloatTok{1}\OperatorTok{,} \VariableTok{n}\NormalTok{)}\OperatorTok{;}
    \VariableTok{X\_bar}\NormalTok{(}\VariableTok{i}\NormalTok{) }\OperatorTok{=} \VariableTok{mean}\NormalTok{(}\VariableTok{sample}\NormalTok{)}\OperatorTok{;}
\KeywordTok{end}

\VariableTok{figure}\OperatorTok{;}
\VariableTok{histogram}\NormalTok{(}\VariableTok{X\_bar}\OperatorTok{,} \FloatTok{60}\OperatorTok{,} \SpecialStringTok{\textquotesingle{}Normalization\textquotesingle{}}\OperatorTok{,} \SpecialStringTok{\textquotesingle{}pdf\textquotesingle{}}\NormalTok{)}\OperatorTok{;}
\VariableTok{hold} \VariableTok{on}\OperatorTok{;}
\VariableTok{x} \OperatorTok{=} \VariableTok{linspace}\NormalTok{(}\FloatTok{0}\OperatorTok{,} \FloatTok{3}\OperatorTok{,} \FloatTok{200}\NormalTok{)}\OperatorTok{;}
\VariableTok{plot}\NormalTok{(}\VariableTok{x}\OperatorTok{,} \VariableTok{normpdf}\NormalTok{(}\VariableTok{x}\OperatorTok{,} \VariableTok{mu\_pop}\OperatorTok{,} \VariableTok{sqrt}\NormalTok{(}\VariableTok{sigma2\_pop}\OperatorTok{/}\VariableTok{n}\NormalTok{))}\OperatorTok{,} \SpecialStringTok{\textquotesingle{}r\textquotesingle{}}\OperatorTok{,} \SpecialStringTok{\textquotesingle{}LineWidth\textquotesingle{}}\OperatorTok{,} \FloatTok{2}\NormalTok{)}\OperatorTok{;}
\VariableTok{xlabel}\NormalTok{(}\SpecialStringTok{\textquotesingle{}Sample Mean\textquotesingle{}}\NormalTok{)}\OperatorTok{;} \VariableTok{ylabel}\NormalTok{(}\SpecialStringTok{\textquotesingle{}PDF\textquotesingle{}}\NormalTok{)}\OperatorTok{;}
\VariableTok{title}\NormalTok{(}\VariableTok{sprintf}\NormalTok{(}\SpecialStringTok{\textquotesingle{}Sampling Distribution of X̄ (n=\%d)\textquotesingle{}}\OperatorTok{,} \VariableTok{n}\NormalTok{))}\OperatorTok{;}
\VariableTok{legend}\NormalTok{(}\SpecialStringTok{\textquotesingle{}Empirical\textquotesingle{}}\OperatorTok{,} \VariableTok{sprintf}\NormalTok{(}\SpecialStringTok{\textquotesingle{}N(\%.1f, \%.4f)\textquotesingle{}}\OperatorTok{,} \VariableTok{mu\_pop}\OperatorTok{,} \VariableTok{sigma2\_pop}\OperatorTok{/}\VariableTok{n}\NormalTok{))}\OperatorTok{;}

\VariableTok{fprintf}\NormalTok{(}\SpecialStringTok{\textquotesingle{}E[X̄] = \%.4f (theory: \%.4f)\textbackslash{}n\textquotesingle{}}\OperatorTok{,} \VariableTok{mean}\NormalTok{(}\VariableTok{X\_bar}\NormalTok{)}\OperatorTok{,} \VariableTok{mu\_pop}\NormalTok{)}\OperatorTok{;}
\VariableTok{fprintf}\NormalTok{(}\SpecialStringTok{\textquotesingle{}Var(X̄) = \%.4f (theory: \%.4f)\textbackslash{}n\textquotesingle{}}\OperatorTok{,} \VariableTok{var}\NormalTok{(}\VariableTok{X\_bar}\NormalTok{)}\OperatorTok{,} \VariableTok{sigma2\_pop}\OperatorTok{/}\VariableTok{n}\NormalTok{)}\OperatorTok{;}
\end{Highlighting}
\end{Shaded}

\end{PSbox}

\begin{PSbox}[unbreakable]{ex}{Example 2: Sample Variance Distribution}

\tcbset{PScodemode/.style={unbreakable}}

\begin{Shaded}
\begin{Highlighting}[]
\CommentTok{\%\% Sampling Distribution of S²}
\CommentTok{\% Population: N(0, 4), σ² = 4}
\VariableTok{sigma2} \OperatorTok{=} \FloatTok{4}\OperatorTok{;} \VariableTok{n} \OperatorTok{=} \FloatTok{10}\OperatorTok{;}
\VariableTok{N\_experiments} \OperatorTok{=} \FloatTok{10000}\OperatorTok{;}

\VariableTok{S2} \OperatorTok{=} \VariableTok{zeros}\NormalTok{(}\FloatTok{1}\OperatorTok{,} \VariableTok{N\_experiments}\NormalTok{)}\OperatorTok{;}
\KeywordTok{for} \VariableTok{i} \OperatorTok{=} \FloatTok{1}\OperatorTok{:}\VariableTok{N\_experiments}
    \VariableTok{sample} \OperatorTok{=} \VariableTok{sqrt}\NormalTok{(}\VariableTok{sigma2}\NormalTok{) }\OperatorTok{*} \VariableTok{randn}\NormalTok{(}\FloatTok{1}\OperatorTok{,} \VariableTok{n}\NormalTok{)}\OperatorTok{;}
    \VariableTok{S2}\NormalTok{(}\VariableTok{i}\NormalTok{) }\OperatorTok{=} \VariableTok{var}\NormalTok{(}\VariableTok{sample}\NormalTok{)}\OperatorTok{;}   \CommentTok{\% Uses n{-}1 denominator}
\KeywordTok{end}

\VariableTok{fprintf}\NormalTok{(}\SpecialStringTok{\textquotesingle{}E[S²] = \%.4f (theory: \%.4f)\textbackslash{}n\textquotesingle{}}\OperatorTok{,} \VariableTok{mean}\NormalTok{(}\VariableTok{S2}\NormalTok{)}\OperatorTok{,} \VariableTok{sigma2}\NormalTok{)}\OperatorTok{;}
\VariableTok{fprintf}\NormalTok{(}\SpecialStringTok{\textquotesingle{}Var(S²) = \%.4f (theory: \%.4f)\textbackslash{}n\textquotesingle{}}\OperatorTok{,} \VariableTok{var}\NormalTok{(}\VariableTok{S2}\NormalTok{)}\OperatorTok{,} \FloatTok{2}\OperatorTok{*}\VariableTok{sigma2}\OperatorTok{\^{}}\FloatTok{2}\OperatorTok{/}\NormalTok{(}\VariableTok{n}\OperatorTok{{-}}\FloatTok{1}\NormalTok{))}\OperatorTok{;}

\VariableTok{figure}\OperatorTok{;}
\VariableTok{histogram}\NormalTok{(}\VariableTok{S2}\OperatorTok{,} \FloatTok{80}\OperatorTok{,} \SpecialStringTok{\textquotesingle{}Normalization\textquotesingle{}}\OperatorTok{,} \SpecialStringTok{\textquotesingle{}pdf\textquotesingle{}}\NormalTok{)}\OperatorTok{;}
\VariableTok{hold} \VariableTok{on}\OperatorTok{;}
\VariableTok{x} \OperatorTok{=} \VariableTok{linspace}\NormalTok{(}\FloatTok{0}\OperatorTok{,} \FloatTok{15}\OperatorTok{,} \FloatTok{200}\NormalTok{)}\OperatorTok{;}
\CommentTok{\% (n{-}1)S²/σ² \textasciitilde{} χ²(n{-}1), so S² \textasciitilde{} σ²/(n{-}1) * χ²(n{-}1)}
\VariableTok{plot}\NormalTok{(}\VariableTok{x}\OperatorTok{,} \VariableTok{chi2pdf}\NormalTok{(}\VariableTok{x}\OperatorTok{*}\NormalTok{(}\VariableTok{n}\OperatorTok{{-}}\FloatTok{1}\NormalTok{)}\OperatorTok{/}\VariableTok{sigma2}\OperatorTok{,} \VariableTok{n}\OperatorTok{{-}}\FloatTok{1}\NormalTok{)}\OperatorTok{*}\NormalTok{(}\VariableTok{n}\OperatorTok{{-}}\FloatTok{1}\NormalTok{)}\OperatorTok{/}\VariableTok{sigma2}\OperatorTok{,} \SpecialStringTok{\textquotesingle{}r\textquotesingle{}}\OperatorTok{,} \SpecialStringTok{\textquotesingle{}LineWidth\textquotesingle{}}\OperatorTok{,} \FloatTok{2}\NormalTok{)}\OperatorTok{;}
\VariableTok{xlabel}\NormalTok{(}\SpecialStringTok{\textquotesingle{}S²\textquotesingle{}}\NormalTok{)}\OperatorTok{;} \VariableTok{ylabel}\NormalTok{(}\SpecialStringTok{\textquotesingle{}PDF\textquotesingle{}}\NormalTok{)}\OperatorTok{;}
\VariableTok{title}\NormalTok{(}\VariableTok{sprintf}\NormalTok{(}\SpecialStringTok{\textquotesingle{}Sampling Distribution of S² (n=\%d, σ²=\%d)\textquotesingle{}}\OperatorTok{,} \VariableTok{n}\OperatorTok{,} \VariableTok{sigma2}\NormalTok{))}\OperatorTok{;}
\VariableTok{legend}\NormalTok{(}\SpecialStringTok{\textquotesingle{}Empirical\textquotesingle{}}\OperatorTok{,} \SpecialStringTok{\textquotesingle{}Theoretical (scaled χ²)\textquotesingle{}}\NormalTok{)}\OperatorTok{;}
\end{Highlighting}
\end{Shaded}

\end{PSbox}

\begin{PSbox}[unbreakable]{ex}{Example 3: Effect of Sample Size}

\tcbset{PScodemode/.style={unbreakable}}

\begin{Shaded}
\begin{Highlighting}[]
\CommentTok{\%\% Precision improves with n: Var(X̄) = σ²/n}
\VariableTok{sigma} \OperatorTok{=} \FloatTok{2}\OperatorTok{;}
\VariableTok{n\_values} \OperatorTok{=}\NormalTok{ [}\FloatTok{5}\OperatorTok{,} \FloatTok{10}\OperatorTok{,} \FloatTok{25}\OperatorTok{,} \FloatTok{50}\OperatorTok{,} \FloatTok{100}\OperatorTok{,} \FloatTok{500}\NormalTok{]}\OperatorTok{;}
\VariableTok{N\_exp} \OperatorTok{=} \FloatTok{10000}\OperatorTok{;}

\VariableTok{fprintf}\NormalTok{(}\SpecialStringTok{\textquotesingle{}n\textbackslash{}t| SE(X̄) Theory\textbackslash{}t| SE(X̄) Sim\textbackslash{}n\textquotesingle{}}\NormalTok{)}\OperatorTok{;}
\VariableTok{fprintf}\NormalTok{(}\SpecialStringTok{\textquotesingle{}{-}{-}{-}{-}{-}{-}{-}{-}+{-}{-}{-}{-}{-}{-}{-}{-}{-}{-}{-}{-}{-}{-}{-}+{-}{-}{-}{-}{-}{-}{-}{-}{-}{-}{-}\textbackslash{}n\textquotesingle{}}\NormalTok{)}\OperatorTok{;}
\KeywordTok{for} \VariableTok{n} \OperatorTok{=} \VariableTok{n\_values}
    \VariableTok{X\_bar} \OperatorTok{=} \VariableTok{mean}\NormalTok{(}\VariableTok{sigma}\OperatorTok{*}\VariableTok{randn}\NormalTok{(}\VariableTok{n}\OperatorTok{,} \VariableTok{N\_exp}\NormalTok{)}\OperatorTok{,} \FloatTok{1}\NormalTok{)}\OperatorTok{;}
    \VariableTok{fprintf}\NormalTok{(}\SpecialStringTok{\textquotesingle{}\%d\textbackslash{}t| \%.4f\textbackslash{}t| \%.4f\textbackslash{}n\textquotesingle{}}\OperatorTok{,} \VariableTok{n}\OperatorTok{,} \VariableTok{sigma}\OperatorTok{/}\VariableTok{sqrt}\NormalTok{(}\VariableTok{n}\NormalTok{)}\OperatorTok{,} \VariableTok{std}\NormalTok{(}\VariableTok{X\_bar}\NormalTok{))}\OperatorTok{;}
\KeywordTok{end}
\end{Highlighting}
\end{Shaded}

\end{PSbox}

\PSsection{12.11}{Practice Problems}

\begin{PSbox}[unbreakable]{prob}{Problem 1}

A population has $\mu = 50$ and $\sigma = 10$. A random sample of $n = 64$ is taken.

\begin{enumerate}
\def\labelenumi{(\alph{enumi})}
\tightlist
\item
  What is $E[\bar{X}]$? (b) What is $\operatorname{Var}(\bar{X})$? (c) What is $P(48 < \bar{X} < 52)$?
\end{enumerate}

\PSsolution{} 

\begin{enumerate}
\def\labelenumi{(\alph{enumi})}
\tightlist
\item
  $E[\bar{X}] = 50$
\item
  $\displaystyle \operatorname{Var}(\bar{X}) = \frac{100}{64} = 1.5625,\, \operatorname{SE} = 1.25$
\item
  $P(48 < \bar{X} < 52) = P(-1.6 < Z < 1.6) = 2\Phi (1.6) - 1 = 0.8904$
\end{enumerate}

\end{PSbox}

\begin{PSbox}[unbreakable]{prob}{Problem 2}

Why $do$ we use $n-1$ in the sample variance formula instead of $n$? Explain both mathematically and intuitively.

\PSsolution{} Mathematically: $E\left[\frac{1}{n}\sum (X_{i}-\bar{X})^{2}\right] = \frac{n-1}{n} \cdot \sigma^{2}$, so dividing by $n-1$ gives $E[S^{2}] = \sigma^{2}$. Intuitively: we use $\bar{X}$ from the same data (not the true $\mu$), so deviations from $\bar{X}$ are artificially small. One degree of freedom is ``used up'' estimating the mean.

\end{PSbox}

\begin{PSbox}[unbreakable]{prob}{Problem 3}

An engineer measures noise power from 20 samples and gets $S^{2} = 0.045\,\mathrm{V}^{2}$. The true variance is $\sigma^{2} = 0.04\,\mathrm{V}^{2}$. Is this concerning? Use the sampling distribution to assess.

\PSsolution{} Under $H_{0}: \sigma^{2} = 0.04,\, \frac{(n-1)S^{2}}{\sigma^{2}} \sim \chi^{2}(19)$. Observed: $\frac{19(0.045)}{0.04} = 21.375$. $P(\chi^{2}(19) > 21.375) \approx 0.31$. This is not unusual — about 31\% of samples would give $S^{2} \ge 0.045$ when $\sigma^{2} = 0.04$. Not concerning.

\emph{Next Module: {Module 13 — Parameter Estimation}}

\end{PSbox}
